\documentclass[11pt]{article}
\usepackage[a4paper,margin=21mm]{geometry}
\usepackage{fontspec}
\setmainfont{Latin Modern Roman}
\newfontfamily\CJKfont{Noto Serif CJK SC}
\XeTeXlinebreaklocale "zh"
\XeTeXlinebreakskip=0pt plus 1pt
\usepackage{amsmath,amssymb,amsthm,amscd,graphicx,color}
\usepackage{xurl}
\usepackage[unicode,hidelinks]{hyperref}
\allowdisplaybreaks
\setlength\emergencystretch{3em}
\numberwithin{equation}{section}

%\usepackage[noabbrev,capitalize]{cleveref}
\usepackage{mathrsfs}
\usepackage{enumitem}
\usepackage{mathtools}

\DeclareMathOperator{\tr}{tr}
\DeclareMathOperator{\dvol}{dvol}
\DeclareMathOperator{\Ric}{Ric}
\DeclareMathOperator{\Rm}{Rm}

\newcommand{\cg}{\widetilde{g}}
\newcommand{\cR}{\widetilde{R}}
\newcommand{\cf}{\widetilde{f}}
\newcommand{\cmG}{\widetilde{\mathcal{G}}}
\newcommand{\cRic}{\widetilde{\Ric}}

\newcommand{\lp}{\langle}
\newcommand{\rp}{\rangle}
\newcommand{\lv}{\lvert}
\newcommand{\rv}{\rvert}

\newcommand{\mF}{\mathcal{F}}
\newcommand{\mG}{\mathcal{G}}
\newcommand{\mW}{\mathcal{W}}

\newcommand{\bN}{\mathbb{N}}
\newcommand{\bR}{\mathbb{R}}
\newcommand{\del}{\delta_{\phi}}
\newcommand{\ric}{\Ric_{\phi}^m}
\newcommand{\psmms}{(M_+,g_+,f_+,m,\mu)}

\numberwithin{equation}{section}

\newtheorem{Theorem}{Theorem}[section]
\newtheorem{Corollary}[Theorem]{Corollary}
\newtheorem{Lemma}[Theorem]{Lemma}
\newtheorem{Proposition}[Theorem]{Proposition}
 { \theoremstyle{definition}
\newtheorem{Definition}[Theorem]{Definition}
\newtheorem{Note}[Theorem]{Note}
\newtheorem{Example}[Theorem]{Example}
\newtheorem{Remark}[Theorem]{Remark} }

\makeatletter\newlabel{ambient-statement}{{1.1}{}{Original source locator}{}{}}\makeatother
\begin{document}
\begin{center}\Large Existence and unique infinite-order jets of straight normal weighted ambient metrics for smooth metric measure spaces with positive noninteger dimension parameter\end{center}
\noindent\textbf{Stable ID:} 1306\quad\textbf{Author:} Jeffrey Case; Ayush Khaitan\par
\noindent\textbf{Work:} The Weighted Ambient Metric\par
\noindent\textbf{Source:} \url{https://sigma-journal.com/2022/086/}\par
\noindent\textbf{Version:} Official SIGMA 18 (2022) article 086, DOI 10.3842/SIGMA.2022.086; journal PDF and native TeX acquired 2026-09-30, exact bytes pinned by SHA256\par
\noindent\textbf{Rights:} CC BY 4.0. Authors retain copyright. \url{https://creativecommons.org/licenses/by/4.0/}. Reuse and typesetting adaptation; original language and mathematical content preserved.\par
\noindent\textbf{Difficulty (prior selection preserved):} {\CJKfont H2: The construction solves coupled tensor and density equations order by order, separating the trace-free metric correction from a two-variable trace/function system. Its determinant identifies both potential resonance orders, while the noninteger dimension hypothesis avoids them. A weighted contracted-Bianchi recurrence then forces the remaining mixed components to vanish to infinite order, completing the formal ambient construction and its jet uniqueness.}\par
\medskip\noindent\textbf{Editorial boundary note (not author text):} Selected Theorem3.8 nonresonant branch: d>=3, finite positive noninteger m, smooth Riemannian g, positive smooth f and real mu. Full printed broader theorem remains context. The outcome is a straight normal dilation-homogeneous Lorentzian metric/function near the cone with weighted Ricci/F vanishing to infinite rho order, unique modulo O(rho\textasciicircum{}infinity), not an exact convergent PDE solution. Complete smooth-measure/curvature/density/weighted-Bianchi and normal/straight/ambient-equivalence definitions are retained. Full current normal-form computation, coupled metric/function correction matrices and determinant, Borel construction4.5(i), remaining-component weighted Bianchi argument and uniqueness are unchanged. Standard unweighted straight-metric characterizations remain precise author references. Resonant/integer warped-product, Poincare, GJMS and example outcomes receive no counts. All source formulas including literal printed curvature notation variants remain unchanged; no editor proof or formula repair.\par
\noindent\textbf{External original reference ambient-statement:} Original unselected ambient-equivalence introductory theorem native129–147; selected straight/normal Theorem3.8 has complete statement/proof here\par
\medskip\hrule\medskip\noindent\textbf{Author text begins below.}\par

% Original source lines 88--104
A \emph{smooth metric measure space} is a five-tuple $(M^d,g,f,m,\mu)$ formed from a Riemannian manifold $(M^d,g)$, a positive function $f\in C^\infty(M)$, a dimensional parameter $m\in[0,\infty]$, and a curvature parameter $\mu\in\bR$. Smooth metric measure spaces arise in many ways, including as (possibly collapsed) limits of sequences of Riemannian manifolds~\cite{CheegerColding2000a}, as smooth manifolds satisfying curvature-dimension inequalities~\cite{BakryEmery1985,Wei_Wylie}, as the geometric framework~\cite{CaseChang2013} for studying curved analogues of the Caffarelli--Silvestre extension for defining the fractional Laplacian~\cite{CaffarelliSilvestre2007}, and, in the limiting case $m=\infty$, as a geometric framework for the realization of the Ricci flow as a gradient flow~\cite{Perelman1}.


A \emph{weighted invariant} is tensor-valued function on smooth metric measure spaces such that $T(N,\Phi^\ast g,\Phi^\ast f,m,\mu) \!=\! \Phi^\ast(T(M,g,f,m,\mu))$ for any smooth metric measure space $(M^d,g,f,m,\mu)$ and any diffeomorphism $\Phi \colon N \mapsto M$.
For example, the \emph{Bakry--\'Emery Ricci curvature} of $(M^d,g,f,m,\mu)$ is
\begin{equation*}
 \Ric_\phi^m := \Ric - \frac{m}{f}\nabla^2f
\end{equation*}
and the \emph{weighted scalar curvature} of $(M^d,g,f,m,\mu)$ is
\begin{equation*}
 R_\phi^m := R - \frac{2m}{f}\Delta f - \frac{m(m-1)}{f^2}\bigl( \lv\nabla f\rv^2 - \mu \bigr) .
\end{equation*}
A smooth metric measure space is often studied as an abstract analogue of the base of the warped product $M^d \times_f F^m(\mu)$ of $(M^d,g)$ with an $m$-dimensional spaceform $F^m(\mu)$ of constant sectional curvature $\mu$, in the sense that when $m$ is a nonnegative integer, most weighted invariants are equivalent to Riemannian invariants of the warped product.
For example, if $m \in \bN_0$, then $\Ric_\phi^m$ is the restriction of the Ricci tensor of $M \times_f F^m(\mu)$ to horizontal (i.e., tangent to $M$) vectors and $R_\phi^m$ is the scalar curvature of $M \times_f F^m(\mu)$.


We say that two smooth metric measure spaces $(M^d,g_i,f_i,m,\mu)$, $i\in\{1,2\}$, are \emph{pointwise conformally equivalent} if there is a function $u\in C^\infty(M)$ such that $g_2={\rm e}^{2u}g_1$ and $f_2={\rm e}^uf_1$.

\setcounter{section}{1}
% Original source lines 229--307
\section{Smooth metric measure spaces}
\label{sec:bg.tex}

We recall some weighted invariants relevant to the geometry of smooth metric measure spaces~\cite{Case2014s}. Let $(M^d,g,f,m,\mu)$ be a smooth metric measure space. We discuss here only the case $m<\infty$ due to the presence of this restriction to our main results. This data determines a volume element
\[
{\rm d}v_\phi := f^m\,\dvol_g .
\]
When $m>0$, we set $\phi:=-m\ln f$, so that
\[
{\rm d}v_\phi = {\rm e}^{-\phi}\,\dvol_g .
\]


In terms of $\phi$, it holds that
\begin{align*}
&\Ric_\phi^m= \Ric + \nabla^2\phi - \frac{1}{m} {\rm d}\phi \otimes {\rm d}\phi,
\\
&R_\phi^m= R + 2\Delta\phi - \frac{m+1}{m}\lv\nabla\phi\rv^2 + m(m-1)\mu {\rm e}^{2\phi/m} .
\end{align*}
Recall that \emph{weighted Schouten tensor $P_\phi^m$} and the \emph{weighted Schouten scalar $J_\phi^m$} of $(M^d,g,f,m,\mu)$ are
\begin{gather*}
 P_\phi^m := \frac{1}{d+m-2}\big(\Ric_\phi^m - J_\phi^m g\big), \qquad
 J_\phi^m := \frac{1}{2(d+m-1)}R_\phi^m .
\end{gather*}
We set
\[
F_\phi^m:= f\Delta f+(m-1)\big(|\nabla f|^2-\mu\big).
\]
Observe that
\begin{equation}\label{fphim-equation}
F_\phi^m=\frac{f^2}{m}\bigl[(d+m-2)\mathrm{tr}_g P_\phi^m -(d+2m-2)J_\phi^m\bigr].
\end{equation}

A \emph{weighted Einstein manifold} is a smooth metric measure space $(M^d,g,f,m,\mu)$ such that
\[
P_\phi^m = \lambda g
\]
for some constant $\lambda\in\bR$.
A weighted Einstein manifold is \emph{quasi-Einstein}~\cite{Case_Shu_Wei} if additionally
\[
J_\phi^m=(d+m)\lambda.
\]
It follows from equation~\eqref{fphim-equation} that this is equivalent to
\[
\Ric_\phi^m=2(d+m-1)\lambda g,\qquad F_\phi^m=-2(d+m-1)\lambda f^2.
\]

We define the \emph{weighted divergence} $\delta_\phi T$ of a tensor field $T \in \Gamma(\otimes^k T^*M)$ by
\[
(\delta_\phi T)(X_1,\dots,X_k):=\sum\limits_{i=1}^d \nabla_{e_i} T(e_i,X_1,\dots,X_k)-T(\nabla\phi,X_1,\dots,X_k),
\]
where $X_1,\dots,X_k\in T_p M$ and $\{e_i\}$ is an orthonormal basis for $T_p M$.

The \emph{weighted Weyl tensor $A_\phi^m$} and the \emph{weighted Cotton tensor ${\rm d}P_\phi^m$} of a smooth metric measure space $(M^d,g,f,m,\mu)$ are
\begin{gather*}
 A_\phi^m(x,y,z,w) := \big(\Rm - P_\phi^m \wedge g\big)(x,y,z,w) ,
 \\
 {\rm d}P_\phi^m(x,y,z) := \nabla_x P_\phi^m(y,z) - \nabla_y P_\phi^m(x,z) ,
\end{gather*}
where $h\wedge k$ denotes the Kulkarni--Nomizu product
\[
(h\wedge k)(x,y,z,w) := h(x,z)k(y,w) + h(y,w)k(x,z) - h(x,w)k(y,z) - h(y,z)k(x,w) .
\]
The \emph{weighted Bach tensor} of $(M^d,g,f,m,\mu)$ is
\[
B_\phi^m := \delta_\phi {\rm d}P_\phi^m - \frac{1}{m}\tr {\rm d}P_\phi^m \otimes {\rm d}\phi + \biggl\lp A_\phi^m, P_\phi^m - \frac{Y_\phi^m}{m}g\biggr\rp ,
\]
where $Y_\phi^m := J_\phi^m - \tr P_\phi^m$ and the contractions are
\begin{gather*}
 (\tr {\rm d}P_\phi^m)(x) := \sum_{i=1}^d {\rm d}P_\phi^m(e_i,x,e_i),
 \\
 \biggl\lp A_\phi^m, P_\phi^m - \frac{Y_\phi^m}{m}g\biggr\rp(x,y) := \sum_{i,j=1}^d A_\phi^m(e_i,x,e_j,y)\biggl(P_\phi^m - \frac{Y_\phi^m}{m}g\biggr)(e_i,e_j).
\end{gather*}

The \emph{weighted Bianchi identity}~\cite{Case2014sd} is
\begin{equation}
\label{intro-weighted-bianchi}
 \del \ric-\frac{1}{2}{\rm d}P_\phi^m-\frac{1}{f^2} F_\phi^m {\rm d}\phi=0.
\end{equation}


% Original source lines 337--464
\section{Weighted ambient spaces}
\label{sec:fg}

Let $(M^d,g,f,m,\mu)$ be a smooth metric measure space with $d\geq 3$ and $m<\infty$. Denote by $\mathcal{E}$ the trivial line bundle over $M$. Let $\mathcal{G}$ be the ray subbundle of $S^2 T^*M\oplus \mathcal{E}$ consisting of all triples $(h,u,x)$ such that $h=s^2 g_x$ and $u=sf(x)$ for some $s\in \bR_+$ and $x\in M$. We define the dilation $\delta_s\colon \mG \mapsto \mG$ by $(h,u,x)\mapsto (s^2 h,su,x)$. Let $T:=\frac{{\rm d}}{{\rm d}s}\delta_s\big|_{s=1}$ denote the infinitesimal generator of dilations. Also, let $\pi \colon (h,u,x) \mapsto x$ be the projection from $\mG$ to $M$. There is a canonical metric measure structure $(\boldsymbol{g},\boldsymbol{f})$ on $\mG$ defined by $\boldsymbol{g}(X,Y)=h(\pi_*X,\pi_*Y)$ and $\boldsymbol{f}(h,u,x)=u(x)$.

Given $(g,f)$, define a coordinate chart $\mG\mapsto \bR_+\times M$, by $(t^2g_x,tf(x),x)\mapsto (t,x)$. Then $T=t\partial_t$.

Consider the embedding $\iota \colon \mG\ \hookrightarrow \mG\times\bR,(t,x)\mapsto (t,x,0)$. Each point in $\mG\times\bR$ can be written as $(t,x,\rho)$. We extend the dilation to $\mathcal{G}\times \bR$ as $\delta_s (t,x,\rho)= (st,x,\rho)$.

We now define weighted pre-ambient spaces and special cases thereof.

\begin{Definition}
$(\cmG,\cg,\cf,m,\mu)$ is called a \emph{weighted pre-ambient space} if
\begin{enumerate}[label=$(\roman*)$]\itemsep=0pt
 \item $\cmG$ is a dilation-invariant neighborhood of $\mG\times \{0\}$;
 \item $\cg$ is a smooth metric of signature $(d+1,1)$ and $\widetilde{f}$ is a smooth function on $\cmG$;
 \item $\delta_s^*\widetilde{g}=s^2 \widetilde{g}$ and $\delta_s^*\widetilde{f}=s\widetilde{f}$; and
 \item $(\iota^*\widetilde{g},\iota^*\widetilde{f})=(\boldsymbol{g},\boldsymbol{f})$.
\end{enumerate}
\end{Definition}

\begin{Remark}
Our results can be extended to pseudo-Riemannian metrics in the usual way (cf.~\cite{FeffermanGraham2012}).
\end{Remark}

\begin{Definition}
A weighted pre-ambient space $(\cmG,\cg,\cf,m,\mu)$ is said to be in \emph{normal form} relative to a metric measure structure $(g,f)$ if
\begin{enumerate}[label=$(\roman*)$]\itemsep=0pt
 \item for each fixed $z\in\mG$, the set of $\rho\in\bR$ such that $(z,\rho)\in \cmG$ is an open interval $I_z$ containing~$0$;
 \item for each $z\in\mG$, the curve on $I_z$ defined as $\rho\mapsto(z,\rho)$ is a geodesic in $\cmG$; and
 \item on $\mG\times \{0\}$, it holds that $\cg=\boldsymbol{g} + 2t\,{\rm d}\rho\,\mathrm{d}t$.
\end{enumerate}
\end{Definition}

\begin{Definition}
A weighted pre-ambient space $(\cmG,\cg,\cf,m,\mu)$ for $(M^d,[g,f],m,\mu)$ is be said to be \textit{straight} if any of the following equivalent properties hold:
\begin{enumerate}[label=$(\roman*)$]\itemsep=0pt
 \item for each $p\in \widetilde{\mathcal{G}}$, the dilation orbit $s\mapsto \delta_s p$ is a geodesic for $\widetilde{g}$;

 \item $\widetilde{g}(2T,\cdot)={\rm d}(\widetilde{g}(T,T))$; or

\item the infinitesimal dilation field $T$ satisfies $\widetilde{\nabla}T=\mathrm{Id}$.
\end{enumerate}
\end{Definition}

Given a point $(t,x,\rho)\in \mathcal{G}\times \bR$, we refer to the $t$ coordinate as the $0$ coordinate, and the $\rho$ coordinate as the $\infty$ coordinate. Coordinates in $M$ are denoted with lowercase latin characters $(i,j,k,\dots)$.

The metric of a normal weighted pre-ambient space takes a special form.

\begin{Lemma}
\label{straight-metric-lemma}
Let $(\cmG,\cg,\cf,m,\mu)$ be a weighted pre-ambient space such that for each $z\in\mG$, the set of all $\rho\in\bR$ such that $(z,\rho)\in\cmG$ is an open interval containing $0$. Then $(\cg,\cf)$ is a \textit{normal} metric measure structure if and only if $\cg_{0\infty}=t$, $\cg_{\infty i}=0$ and $\cg_{\infty\infty}=0$.
\end{Lemma}

\begin{proof}
A normal metric measure structure has
\[
\widetilde{g}_{\infty\infty}\big|_{\rho=0}=\widetilde{g}_{i\infty}\big|_{\rho=0}=0 \qquad\text{and}\qquad \widetilde{g}_{0\infty}\big|_{\rho=0}=t.
\]
Moreover, it also has geodesic $\rho$-lines. The $\rho$-lines are geodesics if and only if the Christoffel symbols $\widetilde{\Gamma}_{\infty \infty I}$, $I\in \{0,i,\infty\}$, vanish. Taking $I=\infty$ gives $\partial_{\rho} \widetilde{g}_{\infty \infty}=0$, and hence $\widetilde{g}_{\infty \infty}=0 .$ Now taking $I=i$ and $I=0$ gives $\widetilde{g}_{\infty i}=0$ and $\widetilde{g}_{\infty 0}=t$.

The converse follows similarly.
\end{proof}
Defining a weighted ambient space required some additional notation.

\begin{Definition}
Let $(\widetilde{S}_{IJ},\widetilde{h})$ be a pair of a symmetric 2-tensor $\widetilde{S}_{IJ}$ and a smooth function $\widetilde{h}$ on an open neighborhood of $\mathcal{G}\times \{0\}\subset\mathcal{G}\times \bR$. For $k,m\geq 0$, we write $(\widetilde{S}_{IJ},\widetilde{h})=O_{IJ}^{k,+}(\rho^m)$ if
\begin{enumerate}[label=$(\roman*)$]\itemsep=0pt
\item $\big(\widetilde{S}_{IJ},\widetilde{h}\big)=O(\rho^{m})$;
\item $\widetilde{S}_{00},\widetilde{S}_{0i}=O(\rho^{m+1})$; and
\item $m f^{-k}\widetilde{h}-g^{ij}\widetilde{S}_{ij}= O(\rho^{m+1})$.
\end{enumerate}
\end{Definition}

We now define a weighted ambient space.

\begin{Definition}
A \emph{weighted ambient space} for $(M^d,g,f,m,\mu)$ is a weighted pre-ambient space $(\cmG,\cg,\cf,m,\mu)$ such that
\begin{enumerate}[label=$(\roman*)$]\itemsep=0pt
\item if $d+m\notin 2\mathbb{N}$, then $\big(\widetilde{\mathrm{Ric}_\phi^m},\widetilde{F_\phi^m}\big)=O_{IJ}(\rho^{\infty})$;

\item if $d+m\in 2\mathbb{N}$, then $\big(\widetilde{\mathrm{Ric}_\phi^m},\widetilde{F_\phi^m}\big)=O^{2,+}_{IJ}\big(\rho^{\frac{d+m}{2}-1}\big).$
 \end{enumerate}
\end{Definition}

A key step in the proof of Theorem~\ref{ambient-statement} is the construction and classification of straight and normal weighted ambient metrics \cite{FeffermanGraham2012}.

\begin{Theorem}
 \label{weighted_ambient_metrics_exist}
 Let $(M^d,g,f,m,\mu)$, $d\geq 3$ and $m<\infty$, be a smooth metric measure space. Then there exists a straight and normal weighted ambient space $(\cmG,\cg,\cf,m,\mu)$ for it. Moreover, if $(\cmG_j,\cg_j,\cf_j,m,\mu)$, $j\in\{1,2\}$, are two such weighted ambient spaces, then
\begin{enumerate}[label=$(\roman*)$]\itemsep=0pt
 \item if $d+m\notin 2\mathbb{N}$, then $\big(\cg_1-\cg_2,-2(\cf_1-\cf_2)\big)=O(\rho^\infty)$; and
 \item if $d+m\in 2\mathbb{N}$, then $\big(\cg_1-\cg_2,-2(\cf_1-\cf_2)\big)=O_{IJ}^{1,+}\big(\rho^{\frac{d+m}{2}-1}\big)$.
 \end{enumerate}
\end{Theorem}

We prove Theorem~\ref{weighted_ambient_metrics_exist} in Section~\ref{formal-theory}.

Note that if $(\cg_1-\cg_2,\cf_1-\cf_2)=O_{IJ}\big(\rho^{\frac{d+m}{2}-1}\big)$, then
\begin{gather*}
{\partial^{\frac{d+m}{2}}_\rho}\big|_{\rho=0}\big(\widetilde{f}_1^m\mathrm{dvol}_{\widetilde{g}_1}- \widetilde{f}_2^m\mathrm{dvol}_{\widetilde{g}_2}\big)
\\ \qquad
{}= \biggl(mf^{-1}{\partial_\rho^{\frac{d+m}{2}}}\big|_{\rho=0}\big(\widetilde{f}_1-\widetilde{f}_2\big)
+\frac{1}{2}g^{ij}{\partial_\rho^{\frac{d+m}{2}}}\big|_{\rho=0} (\widetilde{g}_1-\widetilde{g}_2)_{ij}\biggr)f^m\mathrm{dvol}_{g}.
\end{gather*}
Thus, a geometric interpretation of the improved orders of vanishing in the case that $d+m\in 2\mathbb{N}$ is that the weighted volume element $\widetilde{f}^m\operatorname{dvol}_{\widetilde{g}}$ vanishes to one higher order in $\rho$ than the weighted ambient space $(\widetilde{g},\widetilde{f})$.

The following definition generalizes the notion of uniqueness from Theorem~\ref{weighted_ambient_metrics_exist} to general pre-ambient spaces.

\begin{Definition}
{\sloppy
We say that two weighted pre-ambient spaces $(\widetilde{\mathcal{G}}_{1}, \widetilde{g}_{1},\widetilde{f}_1,m,\mu)$ and $(\widetilde{\mathcal{G}}_{2}, \widetilde{g}_{2}, \widetilde{f}_2,\allowbreak m,\mu)$ for $(M^d,[g,f],m,\mu)$ are \emph{ambient-equivalent} if there exist open sets $\mathcal{U}_{1} \subset \widetilde{\mathcal{G}}_{1}$ and $\mathcal{U}_{2} \subset \widetilde{\mathcal{G}}_{2}$, and a diffeomorphism $\phi\colon \mathcal{U}_{1} \rightarrow \mathcal{U}_{2}$ such that}
\begin{enumerate}[label=$(\roman*)$]\itemsep=0pt
\item $\mathcal{U}_{1}$ and $\mathcal{U}_{2}$ both contain $\mathcal{G} \times\{0\}$;

\item $\mathcal{U}_{1}$ and $\mathcal{U}_{2}$ are dilation-invariant and $\phi$ commutes with dilations;

\item the restriction of $\phi$ to $\mathcal{G} \times\{0\}$ is the identity map; and

\item \begin{itemize}
\item[$(a)$] if $d+m\notin 2\mathbb{N}$, then $\big(\widetilde{g}_{1}-\phi^{*} \widetilde{g}_{2}, -2(\widetilde{f}_1-\phi^*\widetilde{f}_2)\big)=O(\rho^\infty)$;
\item[$(b)$] if $d+m\in 2\mathbb{N}$, then $\big(\widetilde{g}_{1}-\phi^{*} \widetilde{g}_{2}, -2(\widetilde{f}_1-\phi^*\widetilde{f}_2)\big)=O_{I J}^{1,+}\big(\rho^{\frac{d+m}{2}-1}\big)$.
\end{itemize}
\end{enumerate}
\end{Definition}

Using Theorem~\ref{normal-to-straight-and-normal} below, we see that every weighted pre-ambient space is ambient-equivalent to a straight and normal weighted pre-ambient space.



% Original source lines 483--756
\section{Formal theory}\label{formal-theory}

We prove Theorem~\ref{weighted_ambient_metrics_exist} by iteratively constructing a power series solution to $\widetilde{R}(\widetilde{g}),\widetilde{F_\phi^m}=O(\rho^j)$. This process is only obstructed when $d+m\in 2\mathbb{N}$; in this case the obstruction is at order $O\big(\rho^{\frac{d+m}{2}-1}\big)$.

We first give a necessary condition for a weighted ambient space to be normal.

\begin{Lemma}\label{ambient-metric-base-case}
Let $(\widetilde{\mathcal{G}},\widetilde{g},\widetilde{f},m,\mu)$ be a normal weighted ambient space. Then \begin{equation}\label{standard-normal-form}\widetilde{g}=a\,\mathrm{d}t^2+2b_i\, {\rm d}x^i \,\mathrm{d}t+ 2t\,{\rm d}\rho\, \mathrm{d}t+ t^2 (g_{ij})_\rho,\end{equation} with $a=2\rho+O(\rho^2)$ and $b_i=O(\rho^2)$.
\end{Lemma}

\begin{proof}
Lemma~\ref{straight-metric-lemma} implies that $\widetilde{g}$ has the form of equation~\eqref{standard-normal-form}. Since $\widetilde{g}$ is normal, $a|_{\rho=0}=0$ and $b_i|_{\rho=0}=0$. Denote $\widetilde{R}_{IJ}:=\widetilde{(\ric)}_{IJ}$.
By definition,
\begin{gather*}
\begin{split}
&\widetilde{R}_{IJ}=\frac{1}{2}\widetilde{g}^{KL}\big(\partial^2_{IL}\widetilde{g}_{JK} +\partial^2_{JK}\widetilde{g}_{IL}-\partial^2_{KL}\widetilde{g}_{IJ}-\partial^2_{IJ}\widetilde{g}_{KL}\big)
\\ &\hphantom{\widetilde{R}_{IJ}=}
{}+\widetilde{g}^{KL}\widetilde{g}^{PQ}\big(\widetilde{\Gamma}_{ILP}\widetilde{\Gamma}_{JKQ} -\widetilde{\Gamma}_{IJP}\widetilde{\Gamma}_{KLQ}\big)
-\frac{m}{\widetilde{f}}\big(\partial^2_{IJ}\widetilde{f} -\widetilde{\Gamma}^K_{IJ}\partial_K\widetilde{f}\big).
\end{split}
\end{gather*}
Evaluating at $\rho=0$ yields
\begin{gather*}
\widetilde{R}_{00}= \frac{d+m}{2t^2}(2-\partial_{\rho}a),
\qquad
\widetilde{R}_{0i}=\frac{1}{2t}\partial^2_{i\rho}a -\frac{d+m}{2t}\partial_{\rho}b_i.
\end{gather*}
We conclude that $a=2\rho+O(\rho^2)$ and $b_i=O(\rho^2)$.
\end{proof}

We now take a brief digression to discuss straight weighted ambient spaces. First, notice that a simple choice of certain components of $\widetilde{g}$ makes $\widetilde{R}_{0I}$ vanish.

\begin{Lemma}
\label{straight-zero-theorem}
Let $(\widetilde{\mathcal{G}},\widetilde{g},\widetilde{f},m,\mu)$ be a weighted pre-ambient space. If $\widetilde{g}$ has the form
\begin{equation}
\label{straightmetric}
 \widetilde{g}_{IJ}=\begin{pmatrix} 2\rho&0&t\\0& t^2 g_{\rho}&0\\t&0&0 \end{pmatrix}\!,
\end{equation}
where $g_\rho$ is a one-parameter family of metrics on $M$, then $\widetilde{R}_{0I}=0$.
\end{Lemma}

\begin{proof}This follows by direct computation~(cf.\ \cite[Lemma 3.2]{FeffermanGraham2012}).
\end{proof}

The relevance of Lemma~\ref{straight-zero-theorem} stems from the following equivalent characterizations of straight normal pre-ambient metrics.
\begin{Proposition}
\label{straightproperties}
Let $(\widetilde{\mathcal{G}},\widetilde{g},\widetilde{f},m,\mu)$, $d\geq 3$ and $m<\infty$, be in normal form relative to $(g,f)$. Then the following conditions are equivalent:
\begin{enumerate}[label=$(\roman*)$]\itemsep=0pt
\item $\widetilde{g}_{00}=2\rho$ and $\widetilde{g}_{0i}=0$;

\item for each $p\in \widetilde{\mathcal{G}}$, the dilation orbit $s\mapsto \delta_s p$ is a geodesic for $\widetilde{g}$;

\item $\widetilde{g}(2T,\cdot)={\rm d}(\widetilde{g}(T,T))$;

\item the infinitesimal dilation field $T$ satisfies $\widetilde{\nabla}T=\operatorname{Id}$.
\end{enumerate}
\end{Proposition}

\begin{proof}
The proof is identical to that of \cite[Proposition 3.4]{FeffermanGraham2012}.
\end{proof}

The following result implies that, up to ambient-equivalence, we may restrict our attention to metrics of the form of equation~\eqref{straightmetric}.
\begin{Theorem}
\label{normal-to-straight-and-normal}
A normal weighted ambient space is ambient-equivalent to a straight and normal ambient space.
\end{Theorem}
\begin{proof}
Let $(\widetilde{\mathcal{G}},\widetilde{g},\widetilde{f},m,\mu)$ be a normal weighted ambient space. By Lemma~\ref{ambient-metric-base-case}, we may write
\[
\widetilde{g}=\begin{pmatrix} 2\rho&0&t\\0&t^2 g_\rho& 0\\ t&0&0\end{pmatrix}+\rho^n \begin{pmatrix} a& tb_i& 0\\ tb_j & 0 & 0\\ 0&0&0\end{pmatrix}
\]
for some $n\geq 2$ and some $a=a(x,\rho)$ and $b_i=b_i(x,\rho)$. Direct computation using Lemma~\ref{straight-zero-theorem} yields
\begin{gather*}
t^{2} \widetilde{R}_{00}=n\bigg(n-1-\frac{d+m}{2}\bigg) \rho^{n-1} a+O(\rho^{n}),
\\
t \widetilde{R}_{0 i}=n\bigg(n-1-\frac{d+m}{2}\bigg) \rho^{n-1} b_i+\frac{n}{2} \rho^{n-1} \partial_{i} a+O(\rho^{n}),
\end{gather*}
(cf.\ \cite[equation~(3.11)]{FeffermanGraham2012}).
We conclude that if $d+m\notin 2\mathbb{N}$, then $a,b_i=O(\rho^\infty)$; and if $d+m\in 2\mathbb{N}$, then $n\geq \frac{d+m}{2}$. Therefore $\widetilde{g}$ is ambient-equivalent to a metric of the form of equation~\eqref{straightmetric}. The~conclusion follows from Proposition~\ref{straightproperties}.
\end{proof}

We now iteratively construct a straight and normal weighted ambient space.
Let $n\in\mathbb{N}$ be such that there is a metric
\[
\widetilde{g}_{IJ}^{(n-1)} = \begin{pmatrix}
2\rho & 0 & t \\
0 & t^2 g_{\rho} & 0 \\
t & 0 & 0
\end{pmatrix}
\]
and a function $\widetilde{f}^{(n-1)}=tf_\rho$ such that
\begin{gather*}
\widetilde{R}^{(n-1)}=O\big(\rho^{n-1}\big),
\qquad
\widetilde{F_{\phi}^m}^{(n-1)}=O\big(\rho^{n-1}\big).
\end{gather*}
Note that the existence of $\widetilde{g}_{IJ}^{(1)}$ and $\widetilde{f}^{(1)}$ trivially holds. Let $\cg^{(n)}_{IJ}=\cg^{(n-1)}_{IJ}+\Phi_{IJ}$ and $\widetilde{f}^{(n)}= \widetilde{f}^{(n-1)}+\rho^{n}t\upsilon$, where
\begin{align*}
\Phi_{IJ}&= \rho^n \begin{pmatrix}
0 & 0 & 0 \\
0 & t^2\psi_{ij} & 0 \\
0 & 0 & 0
\end{pmatrix}
\end{align*}
and $\psi_{ij}$, $\upsilon$ depend only on $x$ and $\rho$.
We seek $\psi_{IJ}$ and $\upsilon$ such that
 \begin{gather*}
\widetilde{R}^{(n)}= O(\rho^{n}),
\qquad
\widetilde{F_{\phi}^m}^{(n)} = O(\rho^{n}).
\end{gather*}
Note that the inverse of $\cg^{(n)}$, calculated modulo $O(\rho^n)$, is
\[
\cg^{IJ} = \begin{pmatrix} 0 & 0 & t^{-1}\\ 0 & t^{-2}g^{ij}& 0\\ t^{-1} & 0 & -{2\rho}t^{-2}\end{pmatrix}\!.
\]

The Christoffel symbols for $\widetilde{g}^{(n)}$ modulo $O(\rho^n)$ are (cf.~\cite[equation~(3.16)]{FeffermanGraham2012}).
\begin{gather*}
\widetilde{\Gamma}_{I J }^0=\begin{pmatrix}
0 & 0 & 0 \\
0 & -\frac{1}{2} t\big(\partial_{\rho} g_{i j}+n \rho^{n-1} \psi_{i j}\big) & 0 \\
0 & 0 & 0
\end{pmatrix}\!,
\\
\widetilde{\Gamma}_{I J}^\ell=\begin{pmatrix}
0 & t^{-1} \delta_{j}^{\ell} & 0 \\
t^{-1} \delta_{i}^{\ell} & \Gamma_{i j}^\ell & \frac{1}{2} g^{\ell k} \partial_{\rho} g_{i k}+\frac{n}{2} \rho^{n-1} \psi_{i}^{\ell}\\
0 & \frac{1}{2} g^{\ell k} \partial_{\rho} g_{j k}+\frac{n}{2} \rho^{n-1} \psi_{j}^{\ell} & 0
\end{pmatrix}\!,
\\
\widetilde{\Gamma}_{I J}^\infty=\begin{pmatrix}
0 & 0 & t^{-1} \\
0 & \rho \partial_{\rho} g_{i j}-g_{i j} & 0 \\
t^{-1} & 0 & 0
\end{pmatrix}\!.
\end{gather*}
It follows that~(cf.~\cite[equation~(3.11)]{FeffermanGraham2012})
\begin{subequations}
\begin{gather}
\label{ricij}
\widetilde{R}_{ij}^{(n)}=\widetilde{R}_{ij}^{(n-1)}+n\rho^{n-1} \bigg[\bigg(n-\frac{d+m}{2}\bigg)\psi_{ij}-\frac{1}{2}g^{kl}\psi_{kl}g_{ij}-\frac{m}{f}\upsilon g_{ij}\bigg]+O\big(\rho^{n}\big),
\\
\label{rphi}
\widetilde{F_{\phi}^m}^{(n)}= \widetilde{F_{\phi}^m}^{(n-1)}
+ n\rho^{n-1}f \bigg[\frac{1}{2}fg^{kl}\psi_{kl}+\upsilon (d+2m-2n) \bigg]+O\big(\rho^{n}\big),
\\
\widetilde{R}_{\infty\infty}^{(n)}=\widetilde{R}_{\infty\infty}^{(n-1)}-n(n-1)\rho^{n-2} \bigg[\frac{1}{2}g^{kl}\psi_{kl}+mf^{-1}\upsilon\bigg]+O\big(\rho^{n-1}\big)\label{ricinftyinfty}.
\end{gather}
\end{subequations}

\subsection[Solving widetilde {R}\_\{ij\}, widetilde {F\_\{phi\}\textasciicircum{}m}=O(rho\textasciicircum{}n)]{Solving $\boldsymbol{\widetilde{R}_{ij},\widetilde{F_\phi^m}=O(\rho^n)}$}
\label{section-4.1}

The following theorem specifies the result of recursively determining $(\widetilde{g}^{(n)},\widetilde{f}^{(n)})$ by the requirements $\widetilde{R}_{ij}^{(n)},(\widetilde{F}_\phi^m)^{(n)}=O(\rho^n)$.
\begin{Theorem}
\label{various-possibilities}
Let $(M^d,g,f,m,\mu)$, $d\geq 3$ and $m<\infty$, be a smooth metric measure space. Set $\widetilde{\mathcal{G}}:=\mathbb{R}_+\times M\times (-\epsilon,\epsilon)$.
\begin{enumerate}[label=$(\roman*)$]\itemsep=0pt
\item If $d+m\notin\mathbb{N}$, then there is a straight and normal metric measure structure $(\widetilde{g},\widetilde{f})$ on $\widetilde{\mathcal{G}}$, unique modulo $O(\rho^\infty)$, such that $\widetilde{R}_{ij},\widetilde{F_\phi^m}=O(\rho^\infty)$.

\item If $d+m\in 2\mathbb{N}$, then
\begin{enumerate}\itemsep=0pt
\item[$(a)$] there is a straight and normal metric measure structure $(\widetilde{g},\widetilde{f})$ on $\widetilde{\mathcal{G}}$, unique modulo $O^+\big(\rho^{\frac{d+m}{2}}\big)$, such that $(\widetilde{R}_{ij},\widetilde{F_\phi^m})=O_{ij}^{2,+}\big(\rho^{\frac{d+m}{2}-1}\big)$;
\item[$(b)$] if
\[
\partial_\rho^{(d+m)/2-1}\big|_{\rho=0}\widetilde{R}_{ij}=0,
\]
then there is a straight and normal metric measure structure $(\widetilde{g},\widetilde{f})$ on $\widetilde{\mathcal{G}}$ such that $\widetilde{R}_{ij},\widetilde{F_\phi^m}=O(\rho^{d+m-1})$;

\item[$(c)$] if
\begin{equation}\label{oddconsistency}
\frac{m}{f^2} \widetilde{F_\phi^m}^{(d+m-1)}-g^{ij}\widetilde{R}_{ij}^{(d+m-1)}=O\big(\rho^{d+m}\big),
\end{equation}
then there is a straight and normal metric measure structure $(\widetilde{g},\widetilde{f})$ on $\widetilde{\mathcal{G}}$ such that $\widetilde{R}_{ij},\widetilde{F_\phi^m}=O(\rho^\infty)$.
\end{enumerate}

\item If $d+m\in \mathbb{N}\setminus 2\mathbb{N}$, then
\begin{enumerate}\itemsep=0pt
\item[$(a)$] there is a straight and normal metric measure structure $(\widetilde{g},\widetilde{f})$ on $\widetilde{\mathcal{G}}$, unique modulo $O(\rho^{d+m})$, such that $\widetilde{R}_{ij},\widetilde{F_\phi^m}=O(\rho^{d+m-1})$;
\item[$(b)$] if equation~\eqref{oddconsistency} holds, then there is a straight and normal metric measure structure $(\widetilde{g},\widetilde{f})$ on $\widetilde{\mathcal{G}}$ such that $\widetilde{R}_{ij},\widetilde{F_\phi^m}=O(\rho^\infty)$.
\end{enumerate}
\end{enumerate}
\end{Theorem}

\begin{Remark}
We say that $(\widetilde{g},\widetilde{f})$ is unique modulo $O^+(\rho^k)$ if it is unique modulo $O(\rho^{k})$ and $\frac{1}{2}g^{ij}\widetilde{g}_{ij}+\frac{m}{f}\widetilde{f}$ is unique modulo $O(\rho^{k+1})$.
\end{Remark}

\begin{proof}
On studying equations~\eqref{ricij} and \eqref{rphi}, we observe that if $n\neq \frac{d+m}{2}$, then there is a~unique choice of the trace-free part $\psi_{ij}-\frac{1}{d}g^{kl}\psi_{kl}g_{ij}$ of $\psi_{ij}$ which makes the trace-free part of $\widetilde{R}^{(n)}_{ij}$ vanish modulo $O(\rho^n)$. However, if $n=\frac{d+m}{2}$, then we may not be able to make the trace-free part of $\widetilde{R}_{ij}^{(n)}$ vanish modulo $O(\rho^n)$. Additionally, equations~\eqref{ricij} and \eqref{rphi} imply that
\begin{gather}
g^{ij}\widetilde{R}_{ij}^{(n)}=g^{ij}\widetilde{R}_{ij}^{(n-1)}
+n\rho^{n-1}\frac{(2n-2d-m)}{2}g^{kl}\psi_{kl}-\frac{dmn}{f}\upsilon \rho^{n-1}+O\big(\rho^{n}\big),\nonumber
\\
\widetilde{F_{\phi}^m}^{(n)}= \widetilde{F_{\phi}^m}^{(n-1)}
+ n\rho^{n-1}f \bigg[\frac{1}{2}fg^{kl}\psi_{kl}+\upsilon (d+2m-2n) \bigg]+O\big(\rho^{n}\big).
\label{modified-equations}
\end{gather}
The determinant of the coefficient matrix of $g^{kl}\psi_{kl}$ and $\upsilon$ is
\begin{equation*}(2n-d-m)(n-d-m).\end{equation*}
Therefore if $n\notin\{\frac{d+m}{2},d+m\}$, then there are unique $g^{kl}\psi_{kl}$ and $\upsilon$ such that $g^{ij}\widetilde{R}_{ij}^{(n)},(\widetilde{F_\phi^m})^{(n)}=O(\rho^n)$. However, if $n\in\{\frac{d+m}{2},d+m\}$, then we may not be able to simultaneously solve $g^{ij}\widetilde{R}_{ij}^{(n)}=O(\rho^n)$ and $\widetilde{F_\phi^m}=O(\rho^n)$.

Suppose first that $d+m\notin \mathbb{N}$. Combining the above discussion with Borel's lemma yields a~straight and normal metric measure structure $(\widetilde{g},\widetilde{f})$ such that $\widetilde{R}_{ij},\widetilde{F_\phi^m}=O(\rho^\infty)$.

Suppose next that $d+m\in 2\mathbb{N}$. Set $n=\frac{d+m}{2}$. Equations~\eqref{modified-equations} imply that
\begin{align}
\frac{m}{f^2}\widetilde{F_{\phi}^m}^{(\frac{d+m}{2})}- g^{ij}\widetilde{R}_{ij}^{(\frac{d+m}{2})} ={}&\frac{m}{f^2}\widetilde{F_{\phi}^m}^{(\frac{d+m}{2}-1)}-g^{ij}\widetilde{R}_{ij}^{(\frac{d+m}{2}-1)} \nonumber
\\
&+\frac{(d+m)^2}{2}\rho^{\frac{d+m}{2}-1} \bigg(\frac{1}{2}g^{kl}\psi_{kl}+\frac{m}{f}\upsilon\bigg)+O\big(\rho^{\frac{d+m}{2}}\big).
\label{even-difference}
\end{align}
In particular, there is a unique choice of $\frac{1}{2}g^{kl}\psi_{kl}+\frac{m}{f}\upsilon$ such that the left hand side of equation~\eqref{even-difference} is $O(\rho^{\frac{d+m}{2}})$. Making this choice yields a unique straight and normal metric measure structure $(\widetilde{g},\widetilde{f})$ such that $(\widetilde{R}_{ij},\widetilde{F_\phi^m})=O_{ij}^{2,+}(\rho^{\frac{d+m}{2}-1})$.
Moreover, if the weighted obstruction is zero, then $\widetilde{R}_{ij}^{(n)}=O(\rho^{\frac{d+m}{2}})$. Hence, we may iteratively solve equations~\eqref{ricij} and \eqref{rphi} to obtain a straight and normal metric measure structure $(\widetilde{g},\widetilde{f})$ such that $\widetilde{R}_{ij},\widetilde{F_\phi^m}=O(\rho^{d+m-1})$. The possible obstruction is the same as that discussed in the next case.

Suppose finally that $d+m\in \mathbb{N}\setminus 2\mathbb{N}$. Set $n=d+m$. From the discussion of the first paragraph, we may choose the trace-free part of $\psi_{ij}$ such that $\widetilde{R}_{ij}^{(d+m)}$ is pure trace. Additionally, by equations~\eqref{modified-equations},
\begin{gather*}
%\label{odd-consistency-equations}
g^{ij}\widetilde{R}_{ij}^{(d+m)}=g^{ij}\widetilde{R}_{ij}^{(d+m-1)}+
m(d+m) \rho^{d+m-1}\bigg(\frac{1}{2}g^{kl}\psi_{kl}-\frac{d}{f}\upsilon\bigg)+O\big(\rho^{d+m}\big),
\\
\widetilde{F_{\phi}^m}^{(d+m)}= \widetilde{F_{\phi}^m}^{(d+m-1)}
+f^2(d+m)\rho^{d+m-1} \bigg(\frac{1}{2}g^{kl}\psi_{kl}-\frac{d}{f}\upsilon\bigg)+O\big(\rho^{d+m}\big).
\end{gather*}
Therefore we may choose $g^{ij}\psi_{ij}$ and $\upsilon$ such that $g^{ij}\widetilde{R}_{ij}^{(d+m)},\widetilde{F_{\phi}^m}^{(d+m)}=O(\rho^{d+m})$ if and only if
equation~\eqref{oddconsistency} holds. If equation~\eqref{oddconsistency} holds, then we may continue iteratively improving $(\widetilde{g},\widetilde{f})$ as in the first paragraph. Hence, by Borel's lemma, there is a straight and normal metric measure structure $(\widetilde{g},\widetilde{f})$ such that $\widetilde{R}_{ij},\widetilde{F_\phi^m}=O(\rho^\infty)$.
\end{proof}

\subsection[Solving widetilde {R}\_\{I infty\}=0]{Solving $\boldsymbol{\widetilde{R}_{I\infty}=0}$}
In this subsection, we show that if $(\widetilde{g},\widetilde{f})$ is as in Theorem~\ref{various-possibilities}, then the components $\widetilde{R}_{I\infty}$ can be made to vanish to the appropriate order and that equation~\eqref{oddconsistency} automatically holds. This proves Theorem~\ref{weighted_ambient_metrics_exist}.

\begin{proof}[Proof of Theorem~\ref{weighted_ambient_metrics_exist}]
Let $(\widetilde{g},\widetilde{f})$ be as in Theorem~\ref{various-possibilities}. Equation~\eqref{intro-weighted-bianchi} implies that
\begin{gather}
\widetilde{g}^{JK}\widetilde{\partial}_J\widetilde{R}_{KI} -\widetilde{g}^{JK}\widetilde{\Gamma}_{JK}^Q\widetilde{R}_{QI} -\widetilde{g}^{JK}\widetilde{\Gamma}_{JI}^Q\widetilde{R}_{KQ} -\widetilde{g}^{JK}\widetilde{R}_{IJ}\partial_K \widetilde{\phi}
\!-\frac{1}{2}\partial_I \widetilde{R_{\widetilde{\phi}}^m}-\frac{1}{f^2}\widetilde{F}_\phi^m \partial_I\widetilde{\phi}=0.\!
\label{bianchianalogue}
\end{gather}

Set
\begin{equation*}
n=\begin{cases}
\infty & \text{if}\quad d+m\notin\mathbb{N},\\
\frac{d+m}{2}-1&\text{if}\quad d+m\in 2\mathbb{N},\\
d+m-1& \text{if}\quad d+m\in \mathbb{N}\setminus 2\mathbb{N}.
\end{cases}
\end{equation*}
Taking $I=l$ and $I=\infty$ in equation~\eqref{bianchianalogue} and computing$\;\bmod \;O(\rho^{n})$ yields
\begin{subequations}
\label{bianchieqns}
\begin{gather}
\label{bianchi2}
[d+m-2-2\rho\partial_{\rho}] \widetilde{R}_{\infty l}-\rho g^{ks}g'_{sl}\widetilde{R}_{\infty k}+2\rho \widetilde{R}_{l\infty}\partial_\rho \phi+\rho \partial_l\widetilde{R}_{\infty\infty}=O\big(\rho^{n}\big),
\\
[d+m-2-\rho\partial_{\rho}] \widetilde{R}_{\infty\infty}
-\frac{1}{2}\partial_{\rho}\bigg(g^{ij}\widetilde{R}_{ij}-\frac{m}{f^2}\widetilde{F_\phi^m}\bigg)
+g^{ij}\big(\widetilde{\nabla}_\phi\big)_i\widetilde{R}_{j\infty}\nonumber
\\ \qquad
+\rho\big(2\partial_\rho\phi-g^{ij}g'_{ij}\big)\widetilde{R}_{\infty\infty}
=O(\rho^{n})\label{bianchi3}.
\end{gather}
\end{subequations}

First, suppose that $d+m\notin\mathbb{N}$. We know that $\widetilde{R}_{ij},\widetilde{F_\phi^m}=O(\rho^\infty)$. From equations~\eqref{bianchi2} and~\eqref{bianchi3} we conclude that $\widetilde{R}_{\infty l},\widetilde{R}_{\infty\infty}=O(\rho^\infty)$.

Second, suppose that $d+m\in \mathbb{N}\setminus 2\mathbb{N}$. Equation~\eqref{bianchieqns} implies that $\widetilde{R}_{\infty l},\widetilde{R}_{\infty\infty}=O(\rho^{n})$, and that $g^{ij}\widetilde{R}_{ij}-mf^{-2}\widetilde{F_\phi^m}=O(\rho^{n+1})$. This verifies equation~\eqref{oddconsistency}. Equation~\eqref{ricinftyinfty} gives us the unique value of $\frac{1}{2}g^{kl}\psi_{kl}+\frac{m}{f}\upsilon$ such that $\widetilde{R}_{\infty\infty}=O(\rho^{n})$. Hence, we can now uniquely determine the values of $g^{kl}\psi_{kl}$ and $\upsilon$, and solve $\widetilde{R}_{IJ},\widetilde{F_\phi^m}=O(\rho^\infty)$.

Third, suppose that $d+m\in 2\mathbb{N}$. Equation~\eqref{bianchi2} tells us that $\widetilde{R}_{\infty l}=O(\rho^n)$ and that $\widetilde{R}_{\infty\infty}=O(\rho^{n-1})$. Now recall that $g^{ij}\widetilde{R}_{ij}-mf^{-2}\widetilde{F_\phi^m}=O\big(\rho^{\frac{d+m}{2}}\big)$.
Hence, equation~\eqref{bianchi3} tells us that $\widetilde{R}_{\infty\infty}=O(\rho^n)$.

Finally, since the terms of $(\widetilde{g},\widetilde{f})$ are uniquely determined using the process described above, any two straight and normal weighted ambient structures will be ambient-equivalent to the orders stated in the theorem.
\end{proof}
\begin{thebibliography}{99}
\bibitem[2]{BakryEmery1985}
Bakry D., \'Emery M., Diffusions hypercontractives, in S\'eminaire de
 probabilit\'es, {XIX}, 1983/84, \textit{Lecture Notes in Math.}, Vol.~1123,
 \href{https://doi.org/10.1007/BFb0075847}{Springer}, Berlin, 1985, 177--206.

\bibitem[3]{CaffarelliSilvestre2007}
Caffarelli L., Silvestre L., An extension problem related to the fractional
 {L}aplacian, \href{https://doi.org/10.1080/03605300600987306}{\textit{Comm. Partial Differential Equations}} \textbf{32}
 (2007), 1245--1260, \href{https://arxiv.org/abs/math.AP/0608640}{arXiv:math.AP/0608640}.

\bibitem[6]{Case2014sd}
Case J.S., A notion of the weighted {$\sigma_k$}-curvature for manifolds with
 density, \href{https://doi.org/10.1016/j.aim.2016.03.010}{\textit{Adv. Math.}} \textbf{295} (2016), 150--194,
 \href{https://arxiv.org/abs/1409.4455}{arXiv:1409.4455}.

\bibitem[8]{Case2014s}
Case J.S., The weighted {$\sigma_k$}-curvature of a smooth metric measure
 space, \href{https://doi.org/10.2140/pjm.2019.299.339}{\textit{Pacific~J. Math.}} \textbf{299} (2019), 339--399,
 \href{https://arxiv.org/abs/1608.01663}{arXiv:1608.01663}.

\bibitem[10]{CaseChang2013}
Case J.S., Chang S.-Y.A., On fractional {GJMS} operators, \href{https://doi.org/10.1002/cpa.21564}{\textit{Comm. Pure
 Appl. Math.}} \textbf{69} (2016), 1017--1061, \href{https://arxiv.org/abs/1406.1846}{arXiv:1406.1846}.

\bibitem[11]{Case_Shu_Wei}
Case J.S., Shu Y.-J., Wei G., Rigidity of quasi-{E}instein metrics,
 \href{https://doi.org/10.1016/j.difgeo.2010.11.003}{\textit{Differential Geom. Appl.}} \textbf{29} (2011), 93--100,
 \href{https://arxiv.org/abs/0805.3132}{arXiv:0805.3132}.

\bibitem[14]{CheegerColding2000a}
Cheeger J., Colding T.H., On the structure of spaces with {R}icci curvature
 bounded below.~{II}, \href{https://doi.org/10.4310/jdg/1214342145}{\textit{J.~Differential Geom.}} \textbf{54} (2000),
 13--35, \href{https://arxiv.org/abs/1805.07988}{arXiv:1805.07988}.

\bibitem[15]{FeffermanGraham2012}
Fefferman C., Graham C.R., The ambient metric, \textit{Ann. of Math. Stud.},
 Vol.~178, Princeton University Press, Princeton, NJ, 2012.

\bibitem[23]{Perelman1}
Perelman G., The entropy formula for the {R}icci flow and its geometric
 applications, \href{https://arxiv.org/abs/math.DG/0211159}{arXiv:math.DG/0211159}.

\bibitem[24]{Wei_Wylie}
Wei G., Wylie W., Comparison geometry for the {B}akry--{E}mery {R}icci tensor,
 \href{https://doi.org/10.4310/jdg/1261495336}{\textit{J.~Differential Geom.}} \textbf{83} (2009), 377--405,
 \href{https://arxiv.org/abs/0706.1120}{arXiv:0706.1120}.

\end{thebibliography}
\end{document}
